\section{Introduction}

Vector databases~\cite{wang2021milvus, pan2024survey} have moved
from single-tenant prototypes to shared infrastructure that serves
many collections out of the same deployment. A common pattern is
that several collections in the same cluster hold near-duplicate
data: successive versions of a corpus held side by side during a
staged rollout, filtered subsets carved out for specific tenants, or
per-tenant overlays that customize a small fraction of records on
top of a shared substrate. Because these versions overlap heavily
at the raw vector level, a multi-tenant storage layer has an
obvious opportunity to share segments and index structures across
collections rather than duplicating them end to end. Whether that
opportunity is reachable depends on two properties of today's
systems: how much segment-level redundancy actually survives the
ingest path when two near-identical collections are loaded back to
back, and how query throughput behaves once tenants are forced to
share the same physical resources for that redundant data.

This report quantifies both properties on
Milvus~\cite{wang2021milvus}. We frame two layout options for
representing a shared corpus together with per-tenant private
vectors. \emph{Fully replicated} gives each tenant its own collection
and isolates them at every layer of the storage stack at the cost
of duplicating the shared corpus across tenants. \emph{Shared
partition} keeps a single collection in which the shared corpus
and each tenant's private vectors live in distinct partitions, paying
for the shared corpus only once but forcing tenants to contend for
the same memory and page cache at query time. The fully replicated
layout raises the question of whether stock Milvus produces
segment boundaries stable enough for two near-identical collections
to be served from common segments at all. The shared partition
layout raises the question of whether the resulting resource
sharing still delivers competitive query throughput once tenants
run concurrently.

A practical obstacle to studying these questions is that existing
vector search benchmarks target single-tenant accuracy and latency
and do not expose the multi-tenant patterns described above. We
therefore construct two benchmarks on top of established corpora
to probe the two layouts. The first repurposes
CodeSearchNet~\cite{husain2019codesearchnet} from the VIBE
benchmark~\cite{vibe2024} into two near-identical collections that
differ only by a small prefix trim and is loaded under several
Milvus configurations that progressively disable jitter and
compaction, which lets us examine how the storage layer responds
to two tenants whose data is almost identical. The second carves a
Wikipedia corpus into a shared partition and two tenant partitions
under three workload variants that vary how much of the cold
vectors and queries fall on the shared partition, which lets us
examine how query behavior responds when tenants compete for the
same memory and page cache.

The measurements show that segment-level overlap under stock
Milvus is dense and scattered because flush jitter and size-based
compaction push near-identical inputs along divergent merge paths,
and that the underlying workload itself, in which top-100 ground
truth neighbors are spread across many segments per query, leaves
clear room for segment-granularity sharing once that
non-determinism is removed. The shared partition layout, in turn,
trades its storage savings for a sharp throughput cliff under
concurrent tenants, with the signature of page-cache thrashing
rather than CPU saturation. Together these results indicate that
there is real opportunity to optimize the storage layout for
multi-tenant vector databases, and that any such design must
explicitly manage the working set in the page cache rather than
rely on the current segment-per-partition layout and the operating
system's default memory management.
